\section{Fundamentación Teórica}
\subsection{Git and Github}
Git es un sistema de control de versiones distribuido que permite a los desarrolladores gestionar y registrar cada cambio en el código de forma eficiente, segura y con un historial completo de modificaciones. Su arquitectura distribuida posibilita que cada colaborador tenga una copia íntegra del repositorio, lo que incrementa la resiliencia del proyecto y facilita el trabajo sin conexión. GitHub, construido sobre Git, es una plataforma colaborativa en la nube que agrega herramientas para la revisión de código, la gestión de incidencias, las solicitudes de extracción (pull requests) y la integración y entrega continuas. Esta combinación promueve la trazabilidad de cambios, el trabajo organizado en ramas, la colaboración global en proyectos de cualquier escala y la automatización de flujos de desarrollo, contribuyendo a una mayor calidad y velocidad en la entrega de software \cite{ref1}.

\subsection{Scrum}
Scrum es un marco de trabajo ágil para la gestión de proyectos que se fundamenta en ciclos iterativos e incrementales llamados sprints, durante los cuales el equipo planifica, desarrolla y revisa funcionalidades concretas. La versión 2020 de la Guía Scrum refuerza la idea de auto-gestión del equipo, destacando la necesidad de que sus miembros asuman colectivamente la planificación y la responsabilidad de los resultados. También subraya la simplicidad y claridad de sus artefactos principales Product Backlog, Sprint Backlog e Incremento como medio para mantener la transparencia y facilitar el seguimiento del progreso. Los roles esenciales de Scrum Master, Product Owner y Equipo de Desarrollo se definen con mayor precisión para garantizar la comunicación constante y la eliminación de impedimentos. La aplicación de estos principios no solo mejora la adaptabilidad y la capacidad de respuesta ante cambios, sino que también promueve la entrega continua de valor y una cultura de mejora permanente dentro de los proyectos de software \cite{ref2}.

\subsection{Patrones de Diseño}
Los patrones de diseño constituyen soluciones probadas y ampliamente documentadas para resolver problemas recurrentes que surgen durante el desarrollo de software. Se dividen de manera general en tres grandes categorías: creacionales, que se enfocan en la instanciación flexible de objetos; estructurales, que facilitan la composición y relación entre clases y objetos; y de comportamiento, que se centran en la comunicación y asignación de responsabilidades entre componentes. Su propósito principal es favorecer la reutilización de código, la escalabilidad y la mantenibilidad de los sistemas, proporcionando a los desarrolladores un conjunto de guías que ayudan a construir arquitecturas más limpias, coherentes y fáciles de comprender. Además, su adopción fomenta buenas prácticas de programación y acelera la fase de diseño, al ofrecer soluciones previamente validadas por la comunidad de ingeniería de software \cite{ref3}.

\subsection{Modelo-Vista-Controlador (MVC)}
El patrón Modelo–Vista–Controlador (MVC) organiza una aplicación en tres componentes bien definidos. El Modelo se encarga de la lógica de datos y de las reglas de negocio, garantizando la integridad y el acceso a la información. La Vista gestiona la presentación, mostrando al usuario la interfaz y actualizándose conforme cambian los datos. El Controlador actúa como intermediario, recibiendo las acciones del usuario, procesando las solicitudes y dirigiendo el flujo entre el Modelo y la Vista. Esta separación de responsabilidades permite desarrollar, probar y mantener cada capa de manera independiente, favoreciendo la escalabilidad y facilitando la incorporación de nuevas funcionalidades sin afectar el resto del sistema. Además, el MVC mejora la legibilidad del código y promueve la reutilización de componentes, cualidades esenciales en proyectos de mediana y gran envergadura \cite{ref4}.

\subsection{Java}
Java es un lenguaje de programación orientado a objetos ampliamente utilizado en el desarrollo de aplicaciones web, móviles, de escritorio y empresariales. Se distingue por su portabilidad, posibilitada por la Máquina Virtual de Java (JVM), que permite ejecutar el mismo programa en diferentes plataformas sin modificaciones. Además, ofrece características de robustez, seguridad y manejo automático de memoria, lo que reduce errores y mejora la estabilidad de los sistemas. Su extenso ecosistema incluye frameworks y librerías como Spring, Hibernate y JavaFX, que facilitan el desarrollo ágil y escalable, así como herramientas de compilación y automatización que optimizan el ciclo de vida del software. Gracias a su comunidad global activa y a su evolución constante, Java sigue siendo una de las tecnologías preferidas para proyectos de gran envergadura que requieren fiabilidad y rendimiento \cite{ref5}.